\documentclass[a4paper,11pt]{article}

%----------------------------------------------------------------------------------------
%	PACKAGES
%----------------------------------------------------------------------------------------
\usepackage{lmodern}      % vector (Type 1) version of Computer Modern: sharp at any zoom level
\usepackage[T1]{fontenc}
\usepackage[utf8]{inputenc}
\usepackage[french]{babel}
\usepackage{microtype}     % finer spacing and justification
\usepackage{url}
\usepackage[usenames,dvipsnames]{xcolor}
\usepackage[a4paper,top=1cm,bottom=0.8cm,left=1.1cm,right=1.1cm]{geometry}
\usepackage{tabularx}
\usepackage{enumitem}
\usepackage{titlesec}
\usepackage{fontawesome5}
\usepackage[unicode, draft=false]{hyperref}

% links
\definecolor{linkcolour}{rgb}{0,0.2,0.6}
\hypersetup{colorlinks,breaklinks,urlcolor=linkcolour,linkcolor=linkcolour}
\hypersetup{pdftitle={Ihab Jaouad – CV},pdfauthor={Ihab Jaouad},pdfsubject={Stage de fin d'études – Finance de marché, gestion d'actifs, risques, actuariat},pdfkeywords={finance de marché, gestion d'actifs, risques de marché, VaR, Expected Shortfall, P\&L explain, FRTB, Solvabilité II, pricing d'options, Greeks, market making, Python, C++, SQL, VBA}}
% icons are hidden from text extraction (ATS / copy-paste)
\usepackage{accsupp}
\newcommand{\icon}[1]{\BeginAccSupp{ActualText={}}#1\EndAccSupp{}}

% centered version of 'X' column type (header)
\newcolumntype{C}{>{\centering\arraybackslash}X}

% paragraph spacing: controlled explicitly
\setlength{\parindent}{0pt}
\setlength{\parskip}{0pt}

% section style: small caps + rule
\titleformat{\section}{\Large\scshape\raggedright}{}{0em}{}[\titlerule]
\titlespacing{\section}{0pt}{5pt}{3pt}

%----------------------------------------------------------------------------------------
%	CV COMMANDS
%----------------------------------------------------------------------------------------
% \entry{Bold title}{Dates}{Italic subtitle}{Location}
\newcommand{\entry}[4]{%
  \textbf{#1}\hfill #2\par
  \textit{#3}\hfill\textit{#4}\par}
% \entryline{Bold title}{Dates}{Plain description}{Location}
\newcommand{\entryline}[4]{%
  \textbf{#1}\hfill #2\par
  #3\hfill\textit{#4}\par}
% \project{Bold title}{Description}
\newcommand{\project}[2]{%
  \textbf{#1}\par
  #2\par}
% space between two entries
\newcommand{\entrygap}{\par\vspace{3pt}}
% bullet list
\newenvironment{bullets}
  {\begin{itemize}[nosep, leftmargin=1em, itemsep=1pt, label=--]}
  {\end{itemize}}

\begin{document}
\pagestyle{empty}

%----------------------------------------------------------------------------------------
%	EN-TÊTE
%----------------------------------------------------------------------------------------
\begin{tabularx}{\linewidth}{@{} C @{}}
\Huge{Ihab Jaouad} \\[3pt]
{\large Stage de fin d'études (6 mois) – Finance de marché, gestion d'actifs, risques ou actuariat – dès mars 2027}\\
\href{https://www.google.com/maps/place/Paris,+France}{\raisebox{-0.05\height}{\icon{\faMapMarker}}\ Paris – Mobilité : France entière} \ $|$ \
\href{mailto:ihab.jaouad5@gmail.com}{\raisebox{-0.05\height}{\icon{\faEnvelope}} \ ihab.jaouad5@gmail.com} \ $|$ \
\href{tel:+33625586188}{\raisebox{-0.05\height}{\icon{\faMobile}} \ +33 6 25 58 61 88} \\
\end{tabularx}

%----------------------------------------------------------------------------------------
%	PROFIL
%----------------------------------------------------------------------------------------
\section{Profil}
Élève ingénieur à l'ENSEIRB-MATMECA, filière Informatique (data \& IA), classé dans le \textbf{top 3\,\% de la promotion} et en \textbf{double diplôme} avec le \textbf{Master 2 IREF – Finance quantitative et actuariat}. Expérience en modélisation quantitative et data (Python, SQL, C++), des flux de P\&L à une plateforme en production.

%----------------------------------------------------------------------------------------
%	FORMATION
%----------------------------------------------------------------------------------------
\section{Formation}
\entry{Master 2 IREF – Finance quantitative et actuariat}{2026 – 2027}
      {Mathématiques appliquées, statistique – Université de Bordeaux, double diplôme avec l'ENSEIRB}{Bordeaux}
\textit{Cours principaux :} finance mathématique (calcul stochastique), computational finance (Monte Carlo), risque de marché (VaR, Expected Shortfall, backtesting, stress tests), risque de crédit (scoring), réglementation prudentielle (Bâle III), fixed income \& produits dérivés, produits structurés, gestion de portefeuille, actuariat.
\entrygap
\entry{Diplôme d'ingénieur (Bac+5) – Informatique (data \& IA)}{2024 – 2027}
      {ENSEIRB-MATMECA (Bordeaux INP) – \textbf{top 3\,\% de la promotion}}{Bordeaux}
\textit{Cours principaux :} algorithmique numérique, C/C++/Python (POO), IA (ML, deep learning), finance de marché.
\entrygap
\entry{Classes préparatoires aux grandes écoles (CPGE) – MP}{2022 – 2024}
      {Mathématiques et physique intensives – \textbf{major en mathématiques de la promotion}}{Casablanca, Maroc}

%----------------------------------------------------------------------------------------
%	EXPÉRIENCES
%----------------------------------------------------------------------------------------
\section{Expériences professionnelles}
\entry{Quantitative Data Analyst \& Automation – Inside Therapeutics}{juin – sept. 2026}
      {Start-up deeptech – microfluidique}{Bordeaux}
\begin{bullets}
\item Automatisé le traitement de \textbf{$\sim$500 échantillons/mois} avec \textbf{validation indépendante des calculs} : \mbox{réconciliation} (concordance totale), piste d'audit, documentation ; 4\,123 tests (couverture de 91\,\%).
\item Conçu et développé seul LabData de bout en bout, sans cahier des charges : plateforme Python/\allowbreak SQL/\allowbreak Streamlit \textbf{en production en 7 semaines}, faisant gagner \textbf{$\sim$2\,h de travail manuel par expérience} à 6 scientifiques.
\item Pris l'initiative d'ajouter une couche d'analyse (corrélations, régressions, bruit de mesure / variabilité du procédé, alertes d'anomalies) appuyant la \textbf{validation des machines} et l'optimisation des paramètres.
\item Prototypé un modèle par \textbf{processus gaussien} sur plan d'expériences, avec active learning guidé par l'incertitude ; présenté LabData au \textbf{CODIR} ; lancé et animé un workshop IA en interne.
\end{bullets}
\entrygap
\entry{Data Analyst – Finance de marché – Société Générale}{août – sept. 2025}
      {Analyse de données pour plusieurs desks de marché}{Casablanca, Maroc}
\begin{bullets}
\item Intégré \textbf{6 sources de données} (Python, SQL) en flux automatisés et contrôlés alimentant les outils de \textbf{P\&L}.
\item Remplacé des contrôles manuels par \textbf{plus de 25 contrôles automatisés} de qualité des données et de conformité, suivis via des reportings et dashboards Power BI pour les équipes Finance, Risques et Trading.
\end{bullets}
\entrygap
\entryline{Event Manager – Association étudiante CLUBEE}{2024 – 2026}
          {Dirigé une équipe de 5 personnes : plus de 10 événements multiculturels, plus de 100 participants.}{Bordeaux}

%----------------------------------------------------------------------------------------
%	PROJETS
%----------------------------------------------------------------------------------------
\section{Projets}
\project{Pricing, couverture et risques d'options – Python, Streamlit}
{Black–Scholes, Greeks, delta-hedging et \textbf{P\&L explain} ; \textbf{VaR/ES} (FRTB, Solvabilité II), stress tests et backtesting.}
\entrygap
\project{Market making par apprentissage par renforcement – C++20, Python (équipe de 6)}
{Carnet d'ordres C++ (\textbf{400\,k ordres/s}, 2–4\,µs) sur flux Binance ; intégration C++/Python, agent \textbf{Double DQN}.}

%----------------------------------------------------------------------------------------
%	COMPÉTENCES
%----------------------------------------------------------------------------------------
\section{Compétences}
\textbf{Finance \& risques :} pricing d'options (Black–Scholes, Greeks), delta-hedging, P\&L, VaR/Expected Shortfall, backtesting, stress tests, produits dérivés, fixed income, Bâle III/FRTB, Solvabilité II.\par
\textbf{Programmation :} Python (\texttt{pandas}, \texttt{NumPy}, \texttt{scikit-learn}, \texttt{PyTorch}, \texttt{Streamlit}), C/C++, SQL (PostgreSQL), R, VBA ; Java, TypeScript. \textbf{Outils :} Git, Linux, Power BI, Excel, développement assisté par IA (Claude Code).\par
\textbf{Méthodes quantitatives :} calcul stochastique, Monte Carlo, séries temporelles, régression, machine learning.\par
\textbf{Certifications \& distinctions :} certification AMF (en cours), Citi MQA – Quantitative Market Analysis, certification Microsoft Power BI ; lauréat des Olympiades nationales de mathématiques (Maroc).\par
\textbf{Langues :} français (C2), anglais courant (C1 – TOEIC 970/990), arabe (langue maternelle), espagnol (A2).\par
\textbf{Centres d'intérêt :} échecs, football, Formule 1.

\end{document}
